the number of polynomials @@M0@@ is at most equal to the number of connected components of @@M1@@ (EGA III 7.9.5), and it is a small lemma that if, for a @@M2@@, @@M3@@ is locally free of rank @@M4@@, then so is @@M5@@ above an open neighborhood of @@M6@@ in @@M7@@ (see EGA IV 11.3.10, where one takes @@M8@@, @@M9@@). Now it follows from the hypotheses and (EGA III 7.7.6) that, after decomposing @@M10@@ more finely, one can find a sequence @@M11@@ of sections @@M12@@ which is @@M13@@-regular. Thus (ii) follows from the upper semi-continuity property of the function @@M14@@ where, for @@M15@@ and @@M16@@, @@M17@@ denotes the restriction of @@M18@@ to the zero scheme of @@M19@@ (EGA III 7.7.5).

The implication (ii) @@M20@@ (iii) follows immediately from 1.11. The implication (iii) @@M21@@ (iv) follows from 1.3 (iii) by taking @@M22@@ and @@M23@@.

Suppose that @@M24@@ satisfies (iv); there exist exact sequences
@@M25@@
and the set of polynomials @@M26@@ is finite. By hypothesis, the family of classes of the @@M27@@ is bounded; thus, by the implication (i) @@M28@@ (ii) already proved, the set of polynomials @@M29@@ is finite and there exists a sequence of integers @@M30@@ such that all the @@M31@@ (with @@M32@@ algebraically closed) are @@M33@@-sheaves. Consequently, the set of polynomials @@M34@@ is finite and, by 1.6, all the @@M35@@ are @@M36@@-sheaves. Thus, applying to the family of classes of the @@M37@@ the implication (ii) @@M38@@ (iv) already proved, we obtain (v).

Finally, suppose that @@M39@@ satisfies (v); let us prove that @@M40@@ is bounded. One may suppose @@M41@@ (resp. @@M42@@) of the form @@M43@@. Indeed,
